
\chapter{Methodology}
\label{ch:methodology}

This chapter describes how the experiments were run. One shared pipeline evaluates small language
models across memory strategies and model sizes on a smartphone and on a workstation. The evaluation
harness, question sets, configuration baselines and quantization tier are the same in both settings;
the two platforms differ in their hardware, runtime software and weight packaging.

\paragraph{Contribution of this chapter.} The individual components are established: dense turn
retrieval~\citep{dpr}, a reference-guided automatic judge using the LongMemEval
rubric~\citep{longmemeval}, Coulomb counting of battery charge, and accuracy per joule as an
efficiency measure~\citep{intelligence_per_watt}. The contribution is the protocol that combines
them for conversational memory on a phone. Two of its elements distinguish it from the on-device
studies reviewed in Section~\ref{sec:ch2_device}. First, every compared memory strategy runs on the
phone, and its accuracy and energy come from the same phone workload: the answers it produced there
are the ones judged, and the energy is measured over the span that produced them. In the reviewed
studies, either the competing systems did not run on the phone~\citep{smartrag} or the workload was
document retrieval rather than a conversation history~\citep{mobilerag_energy}. Second, a rule fixed
before the comparison decides whether workstation accuracy may stand in for phone accuracy. A third
element serves the graph analysis of Chapter~\ref{ch:graph_mechanisms}: the
graph pipeline is changed one stage at a time (Table~\ref{tab:graph_versions}), so that losses in
retrieval, extraction and representation can be told apart.

\section{Research design: one pipeline, two measurement settings}
\label{sec:design}

The study compares three ways of giving a small language model access to a long conversation, at
several model sizes, on two long-context benchmarks. Running every condition on the phone is not
practical: the full grid spans hundreds of paired questions across five readers, two benchmarks and
a sweep of retrieval depths, while one 20-question full-context sitting on the phone took up to about
four hours. Energy, battery drain, thermal throttling and memory pressure, however, exist only on the
real device. The two settings therefore have complementary jobs:

\begin{itemize}
  \item The phone (a Samsung Galaxy S23 Ultra, running models locally through the
        llama.cpp runtime~\citep{llamacpp}) is the main measurement platform. It measures
        whole-device energy, battery charge, latency, memory use and heat, and it records the
        answers it generates, which the same judge then scores under the same rubric. Phone
        accuracy is therefore measured on the phone, not assumed from the workstation.
  \item The workstation (referred to as ``the box'') runs the same memory modules through the
        Ollama runtime~\citep{ollama} on more questions: up to 200 LoCoMo and 60 LongMemEval-S
        questions per condition, five readers, and the graph and retrieval-depth variants. It also
        hosts the 70B judge. Its larger samples allow paired statistical tests of the main strategy
        comparisons.
\end{itemize}

Phone accuracy and energy per query are reported separately, and correct answers per unit of energy
combines the answers and energy of the same phone workload. Because the runtimes and weight packaging
differ, platform equivalence was tested rather than assumed: a predefined parity rule
(Section~\ref{sec:validity}) decided whether workstation accuracy could substitute for phone accuracy,
and every on-device claim uses phone-measured accuracy. The one exploratory projection that combines
the two platforms is labelled as such in Appendix~\ref{app:ch5_routing_sensitivity}.

The comparison is \emph{paired} throughout: for a given benchmark and model, all memory strategies see
the same questions, dialogue histories and reference answers, which enables paired tests and guards
against data-handling mistakes. Figure~\ref{fig:ch3_architecture} gives an overview of the pipeline.

\begin{figure}[htbp]
  \centering
  \includegraphics[width=\textwidth]{results/thesis_visuals/fig13_architecture_diagram.pdf}
  \caption{Shared memory and evaluation pipeline. Upper rows: truncated full context, vector retrieval and
  SLM-extracted graph memory construct evidence for the same reader and answer prompt; GraphRAG v4
  adds source turns to v3 triples. (B) The workstation evaluates all five readers through
  Ollama. (C) The Galaxy S23 Ultra runs selected configurations through llama.cpp and records
  device resources. A workstation-hosted 70B judge scores both platforms' answers; their
  accuracies remain separate.}
  \label{fig:ch3_architecture}
\end{figure}

\section{Language models and runtime environment}
\label{sec:models}

The model that reads the supplied conversational context and produces the answer is called the
\emph{reader}. Five readers are evaluated, drawn from two open-weight families and spanning from
0.8 to 4 billion parameters (Table~\ref{tab:readers}). Using two families shows whether the
comparison between memory strategies holds beyond one model family.

\begin{table}[t]
  \centering
  \begin{tabular}{llll}
    \toprule
    Family & Reader & Parameter scale & Attention design \\
    \midrule
    Gemma 4~\citep{gemma4}   & E2B  & 2.3B effective (5B total) & Hybrid local/global attention \\
    Gemma 4~\citep{gemma4}   & E4B  & 4.5B effective (8B total) & Hybrid local/global attention \\
    Qwen3.5~\citep{qwen35}   & 0.8B & 0.8B         & Hybrid linear/full attention \\
    Qwen3.5~\citep{qwen35}   & 2B   & 2B           & Hybrid linear/full attention \\
    Qwen3.5~\citep{qwen35}   & 4B   & 4B           & Hybrid linear/full attention \\
    \bottomrule
  \end{tabular}
  \caption{The five reader models evaluated across the memory strategies.}
  \label{tab:readers}
\end{table}

\paragraph{Model architectures.} The two families differ in how they handle long inputs. The
Gemma~4 readers use a hybrid attention scheme that interleaves local sliding-window attention with
full global attention, with the final layer always global; the global layers apply proportional
rotary position encoding, and many layers share the key--value cache of earlier layers, both of which reduce the
memory cost of long contexts~\citep{gemma4}. The ``E'' in E2B and E4B stands for \emph{effective} parameters: the smaller
Gemma models use per-layer embeddings, giving each decoder layer its own small lookup table per
token, so the effective parameter count is much lower than the raw total. These design choices matter
here because they bear directly on the on-device memory analysis presented later. The Qwen3.5 readers
are dense models (no mixture of experts) with a different hybrid design: in every block of four
layers, three use Gated DeltaNet linear attention and one uses gated full attention~\citep{qwen35}.
Gemma~4 is the primary reader family of this thesis; the Qwen3.5 readers were added as a second
family to test whether the findings hold beyond Gemma.

\paragraph{Quantization and packaging.} Running local inference on mobile hardware requires weight
compression. All five readers are evaluated at one quantization tier, the standard \texttt{Q4\_K\_M}
format, which stores weights at roughly four bits each with a fixed per-layer layout. The phone runs
the Unsloth-distributed GGUF files directly through llama.cpp~\citep{unsloth_gguf,llamacpp}, while
the workstation serves the same \texttt{Q4\_K\_M} tier through Ollama~\citep{ollama}. These are the
standard static \texttt{Q4\_K\_M} builds, not the adaptive-width Dynamic variants. One on-device
cell was repeated on an adaptive 4-bit build (\texttt{UD-Q4\_K\_XL}) to check whether the choice of
4-bit build changes the results. It returned the same verdict on all twenty questions as the
standard build, so one build is used throughout. Higher-precision weights were not tested.

\paragraph{On-device hardware and runtime.} On-device measurements were performed on a Samsung
Galaxy S23 Ultra. Inference ran on the CPU of its Snapdragon 8 Gen 2 (one Cortex-X3 prime core at
3.36~GHz, four performance cores at 2.80~GHz and three efficiency cores at 2.00~GHz), with 12~GB of
RAM and a 5,000~mAh battery; temperature was read from the phone's thermal zones. Inference ran
through a local inference server inside a Termux environment, with
execution pinned to the five high-performance cores, five worker threads, and FlashAttention
enabled. The phone server ran with an 8,192-token context. The llama.cpp server rejects a prompt
longer than its context, and full-context prompts built to the 4,096-token budget are longer than
4,096 tokens (Section~\ref{sec:strategies}), so the server was given this room to keep them from
failing. The workstation was a dedicated Linux host running Ollama on GPU hardware. Ollama does not
reject an over-long prompt; it trims it to its 4,096-token window.

\paragraph{Decoding.} Every reader generates its answer by greedy decoding: taking the most likely
token at each step, with temperature set to zero ($T=0.0$). Greedy decoding removes sampling
randomness. On the phone, repeated sittings returned identical answers on every question they
shared, so the variation between repeated sittings reflects physical energy and runtime differences,
not differences in what the model said. Model cards typically document a non-zero sampling
temperature (for example $T=1.0$), so greedy decoding is an intentional deviation, chosen because
deterministic decoding is standard for benchmarking and makes repeated runs directly comparable; the
paired tests themselves rest on every condition answering the same questions. A
control experiment on LoCoMo found no detectable difference between greedy decoding and the vendor
sampling setting (36.0\% against 35.0\%, $p=0.69$, McNemar's exact test).

\paragraph{Step-by-step reasoning.} Some readers can write out intermediate reasoning before the
final answer. This is toggled through the vendor-documented request parameter. In the baseline
configuration reasoning is off, and the models return concise answers of about a dozen tokens. In the
dedicated reasoning experiments reasoning is on and the generation budget is raised from 64 to 768
tokens.

\section{Memory strategies}
\label{sec:strategies}

Three ways of giving the reader access to past dialogue turns are compared. The reader, the decoding
settings, and the answer format are held constant; only the way the context is built differs.
All three strategies work from the same rendered turns. When a LoCoMo turn shares a photo, the
caption that the dataset provides for that photo is appended to the turn's text.

\begin{description}
  \item[Full context (concatenation).] The whole conversation is placed in the prompt in
        chronological order, each turn rendered with its timestamp, speaker, and text. The oldest
        turns are dropped until the history fits a 4,096-token budget, less 192 tokens reserved for
        the template and the answer. The harness counts tokens as characters divided by four, which
        undercounts LoCoMo dialogue, so the resulting prompts measured 4,670 to 5,254 tokens. The
        phone server held them in full. Ollama on the workstation trimmed the oldest part of each
        prompt to exactly 4,096 tokens, so workstation full context read somewhat less history than
        phone full context. On the twenty questions that both platforms answered, the wording of the
        answers differed, but the judged verdicts agreed on 18 or 19 of 20 for every reader. An
        exploratory workstation rerun for three readers raised the window to 131,072 tokens
        (262,144 for one cell), so that the history fit whole in every prompt except four
        LongMemEval-S prompts per Gemma reader. Those four lost their oldest 2 to 4 turns
        (Section~\ref{sec:ch4_long_fc}).
  \item[Vector retrieval (dense turn retrieval).] Each turn is embedded once at ingestion time. At
        query time the question is embedded, its cosine similarity to every stored turn is computed,
        and the most similar turns (by default the top ten) are retrieved, sorted back into
        chronological order, and placed in the prompt.
  \item[Graph memory (SLM-extracted graph).] The reader model itself acts as the extractor while the
        conversation is ingested. Reading the dialogue in batches of five turns, it writes out
        factual relationships as (subject, relation, object) triples in structured form. The facts
        are stored in a directed multigraph, from which the facts relevant to a question are
        retrieved at query time. Because the graph is built by the small model itself rather than by
        an external rule-based parser or a larger model, it is called an \emph{SLM-extracted graph}
        throughout.
\end{description}

Figure~\ref{fig:ch3_lookup_methods} sets the three lookups side by side: what each strategy stores
when the conversation is ingested, what it does when a question arrives, and what text the reader
finally receives.

\begin{figure}[htbp]
  \centering
  \includegraphics[width=\textwidth]{results/thesis_visuals/fig22_lookup_methods.pdf}
  \caption{How each memory strategy looks up evidence (schematic; the turns are illustrative).
  Full context performs no search: it keeps the newest turns that fit the 4,096-token budget, less
  192 tokens, and drops the oldest. Vector retrieval scores every stored turn embedding against the
  question embedding and passes the ten most similar turns to the reader in time order. Graph memory
  (GraphRAG v2--v4) uses the same top ten turns to find edges extracted from their five-turn batches,
  ranks those edges by the similarity of their best turn multiplied by the edge weight, and keeps
  twelve; v3 shows them as triples and v4 adds each source turn once. GraphRAG v1 instead matched
  question terms to node names and took the one-hop neighbourhood (Table~\ref{tab:graph_versions}).}
  \label{fig:ch3_lookup_methods}
\end{figure}

\subsection{Vector retrieval}
\label{sec:vector_details}

Vector retrieval is a standard retrieval-augmented generation
arrangement~\citep{rag_lewis} built on a dual-encoder dense retriever~\citep{dpr}. There is no
reranking stage, no query rewriting or expansion, and no summarization or compression of the
retrieved text. The design choices made here are the unit of retrieval and the decision to add
nothing on top of the retriever.

\paragraph{Unit and ingestion.} The retrieval unit is a single dialogue turn, rendered with its
timestamp, speaker, and text exactly as the full-context condition renders it. Each turn is embedded
once when it is ingested, and turns longer than roughly 1500 tokens are truncated before embedding.
Nothing is re-embedded at query time and nothing is embedded twice, so one ingestion pass serves
every question asked of that conversation and every retrieval depth swept over it.

\paragraph{Query time.} The question is embedded with the query-side prefix, cosine similarity to
every stored turn is computed as a dot product over unit-normalized vectors, and the $k$ highest
scoring turns are taken. Those turns are then sorted back into chronological order before they enter
the reader prompt, so the reader sees dialogue in the order it happened, not in similarity
rank. The default depth is $k = 10$; Chapter~\ref{ch:memory_evaluation} sweeps it from 2 to 50.

\paragraph{Embedder.} Nomic Embed Text v1.5~\citep{nomic_embed} is used off the shelf, with no
fine-tuning on either benchmark and no training or tuning for this study. It is a
general-purpose embedder trained on large collections of text pairs, and general-purpose dense
retrievers are known to work on new collections without task-specific training~\citep{contriever}.
The retrieval arm therefore has no advantage from tuning to these benchmarks. The task prefixes this model expects, and the one respect in which the workstation
and on-device runs differ in applying them, are described in Section~\ref{sec:retrieval}.

\paragraph{What is stored.} The ingestion artifact is a set of per-turn embeddings that depends only
on the embedding model and on whether the task prefixes were applied, not on the reader or the
retrieval depth; depth is applied at query time by slicing the ranked list. Every reader compared at
a given depth therefore sees byte-identical retrieved turns, and every depth in a sweep reads the same
stored embeddings, so differences across readers or depths reflect what was shown, not how it was
indexed.

\subsection{Graph construction and retrieval}
\label{sec:graph_details}

The base SLM-extracted graph works as follows.

\paragraph{Extraction and normalization.} Turns are ingested in non-overlapping, consecutive batches
of five. The extractor is prompted to identify factual relationships, to preserve dates, times, and
locations, to resolve pronouns to explicit speaker names where the reference is clear, and to return a
list of triples. Extraction is limited to 256 tokens of output in the baseline (raised to 512 or 1024 in
specific repair runs). The extracted text is stripped of formatting markers and parsed; entity and
relation strings are then normalized deterministically by lower-casing, removing trailing
punctuation and possessives, and collapsing familial aliases to a canonical form (for example,
\emph{mom} to \emph{mother}). A short stop-list of bare pronouns filters out ungrounded nodes.

\paragraph{Graph structure and reinforcement.} Facts are stored as a directed multigraph. A new
triple adds a weighted edge between its two entities, carrying the dialogue timestamp and the source
turn range; a repeat of the same relation between the same pair of entities reinforces the existing
edge's weight up to a ceiling and updates its timestamp and sources. The implementation supports edge
decay and pruning, but decay is disabled in every reported run (the decay rate is fixed so weights
never shrink), so the evaluated graph is strictly append-only with reinforcement. No claim about
forgetting is made.

The four evaluated versions, GraphRAG v1 to v4, form a controlled improvement ladder
(Table~\ref{tab:graph_versions}). They are named after the whole retrieval pipeline, not the graph:
every version stores the same kind of triple edges in a directed multigraph of the same form.
Each version changes one part of the preceding configuration, which separates a retrieval problem from
an extraction-budget problem or a loss of evidence during graph compression; the v1-to-v2 step also
adds source links and name standardization (Chapter~\ref{ch:graph_mechanisms}).

\begin{table}[t]
  \centering
  \small
  \begin{tabular}{@{}p{0.08\linewidth}p{0.30\linewidth}p{0.12\linewidth}p{0.19\linewidth}p{0.19\linewidth}@{}}
    \toprule
    Version & Retrieval mechanism & Extraction budget & Graph contents & Context given to reader \\
    \midrule
    v1 & Fuzzy entity matching and one-hop ego-subgraph & 256 tokens & Original build & Retrieved triples \\
    v2 & Dense turn search followed by associated-triple ranking & 256 tokens & Rebuilt with source links and standardized names & Retrieved triples \\
    v3 & Same vector-first retrieval as v2 & 512 tokens & Rebuilt; more triples survive extraction & Retrieved triples \\
    v4 & Same retrieval as v3 & 512 tokens & Same as v3 & Triples and verbatim source turns \\
    \bottomrule
  \end{tabular}
  \caption{Configuration of GraphRAG v1 to v4. The triple schema is the same in every version; the
  columns show which stage each version changes.}
  \label{tab:graph_versions}
\end{table}

The v1-to-v2 step asks whether entity matching retrieves the wrong part of the conversation, v2-to-v3
whether structured output cut off at 256 tokens loses facts, and v3-to-v4 whether triples preserve
enough evidence or the original wording is still needed. Version~4 reuses the v3 graph and adds the
supporting source turns when the answer prompt is assembled.

\section{Benchmarks and task configuration}
\label{sec:benchmarks}

Two conversational long-context benchmarks are used.

\paragraph{LoCoMo}~\citep{locomo} contains long, multi-session dialogues between recurring speakers,
with paired questions. The benchmark defines several tasks; this study evaluates only the
\emph{question-answering} task, and every result statement refers to that task. Evaluation uses a
fixed set of 200 questions, 50 from each of Categories 1 through 4, taken in turn from each of the
ten conversations (round-robin) so that every conversation contributes. The benchmark paper reports
a fifth, adversarial category as well~\citep{locomo}; later work that evaluates on LoCoMo commonly
reports only the first four, and this thesis follows that convention. \citet{budgetmem}, for example,
state the same exclusion when reporting per-category LoCoMo results. The exclusion keeps the
results comparable with that work and is not a judgement about the omitted category.

\paragraph{LongMemEval-S}~\citep{longmemeval} provides much longer histories, averaging about 500
dialogue turns, with each question labelled by one of six reasoning types. Because the questions are
stored grouped by type, a simple prefix of the file would sample only one type; to avoid that, a
stratified subset of 60 questions is used, ten from each of the six types, drawn with a fixed random
seed so the subset is reproducible.

\paragraph{Roles of the two benchmarks.} LoCoMo is the primary benchmark: it carries the full
experimental grid, the graph-repair and selective-expansion studies, the answer-prompt substitution
and all on-device measurements. LongMemEval-S tests whether the main strategy ordering holds at
500-turn scale; it carries the core memory-strategy comparison, a retrieval-depth sweep and a
reasoning probe. Where a result exists for only one benchmark, the text says so.

\paragraph{Balanced sampling and population weighting.} Both benchmarks are sampled with equal
questions per category, which gives every category the same number of questions for fine-grained
error analysis, so the plain sample mean is an unweighted category average. Because the real benchmark
populations are not evenly split across categories (for example, Category~4 is about 55\% of the four
evaluated LoCoMo categories), a population-weighted accuracy is also computed as a secondary
sensitivity check.
Table~\ref{tab:ch3_benchmarks} summarizes both evaluated datasets, and
Figure~\ref{fig:ch3_benchmark_anatomy} shows the structure of one LoCoMo history and where the
evidence for one question lies.

\begin{table}[htbp]
  \centering
  \small
  \begin{tabular}{@{}p{0.22\linewidth}p{0.34\linewidth}p{0.34\linewidth}@{}}
    \toprule
     & LoCoMo & LongMemEval-S \\
    \midrule
    Histories & 10 conversations & 60 histories, one per question \\
    Turns per history & 369--689 (mean 588) & 396--561 (mean 491) \\
    Sessions per history & 19--32 & 39--61 (mean 48) \\
    Words per history & 8.0k--16.2k (mean 13.4k) & 74k--82k (mean 79k) \\
    Questions evaluated & 200, 50 from each of Categories 1--4 & 60, 10 from each of six types \\
    Phone runs & a 20-question slice & not run on the phone \\
    Role & primary benchmark & tests whether the ordering transfers \\
    \bottomrule
  \end{tabular}
  \caption{The two evaluated datasets. Sizes are computed from the evaluated data files; token
  counts depend on the reader's tokenizer and are not shown.}
  \label{tab:ch3_benchmarks}
\end{table}

\begin{figure}[htbp]
  \centering
  \includegraphics[width=\textwidth]{results/thesis_visuals/fig00_benchmark_anatomy.pdf}
  \caption{One LoCoMo history, \texttt{conv-26}: 419 turns across 19 sessions. The labelled evidence
  for ``What did Caroline research?'' lies in session 2, outside the 79 newest turns that full context
  retains. Those 79 turns occupy 4,762 tokens on the phone; the workstation trims the prompt further
  to 4,096. The strip below the sessions draws all 419 turns to scale.}
  \label{fig:ch3_benchmark_anatomy}
\end{figure}

\subsection{Fixed configuration and the one-variable rule}
\label{sec:config}

Apart from the variable under study, every run uses one fixed configuration. Most comparisons change a
single element of this baseline. Three change more than one and are reported as combined changes:
GraphRAG v2 (retrieval, source links and name standardization together), the reasoning runs (reasoning
mode together with a longer answer limit), and the graph model-size comparison (extractor and reader
together).

\begin{itemize}
  \item \textbf{History budget.} 4,096 tokens across all models and conditions, which sets where
        full context truncates the history; the only exception is the long-window full-context
        rerun of Section~\ref{sec:ch4_long_fc}. Section~\ref{sec:strategies} describes how each platform
        applies it.
  \item \textbf{Decoding.} Greedy ($T=0.0$), reasoning off.
  \item \textbf{Answer length.} A generation budget of 64 tokens. In practice this ceiling is almost
        never reached: the reader is asked for a short, direct answer and returns about ten to
        fifteen tokens. The reasoning runs raise the budget to 768 tokens, because the reasoning has
        to be written out; that change is confined to those runs and is treated as a controlled
        variable, not part of the fixed baseline.
  \item \textbf{Reader prompt.} One curated short-answer prompt is given to the reader across all
        models, strategies, and benchmarks.
  \item \textbf{Retrieval.} Ten retrieved turns by default, ranked by cosine similarity using Nomic
        Embed Text v1.5 with its task prefixes (Section~\ref{sec:retrieval}). Turns are truncated to
        about 1500 tokens before embedding.
  \item \textbf{Graph construction.} Extractor batches of five turns, a 256-token extraction budget,
        up to twelve retrieved facts per question, decay disabled, and the extractor model set to the
        same model that answers.
  \item \textbf{Scoring.} The same 70B judge and the same judge prompt for every answer, on both
        benchmarks and both devices (Section~\ref{sec:judge}).
\end{itemize}

The varied elements are the memory strategy, the model size, the number of retrieved turns, whether
reasoning is on, the graph's extraction budget and the answer prompt, with the three combined changes
named above. Everything else stays as above.

\section{Retrieval configuration}
\label{sec:retrieval}

Dense retrieval and vector-first graph retrieval use Nomic Embed Text
v1.5~\citep{nomic_embed,nomic_embed_card}, which expects asymmetric task prefixes:
\texttt{search\_document:} for stored turns and triple strings, and \texttt{search\_query:} for
questions. The prefixes were added partway through the study, and the affected workstation
conditions were rerun (Appendix~\ref{app:reproducibility}).

\paragraph{Reporting convention.} A completed prefixed rerun replaces the earlier unprefixed result
for the same condition, whether its score is higher or lower; historical values appear only where
needed to explain development or fixed evaluation criteria, and are identified as such. Three
configurations result:
\begin{itemize}
  \item all primary workstation comparisons use prefixed retrieval;
  \item the primary 20-question phone comparison keeps its original unprefixed runs, because the
        replicated campaign was too costly to repeat, and footnotes mark this exception;
  \item the E2B and E4B 100-question phone runs and the E2B two-turn phone probe use prefixes.
\end{itemize}
Full context and GraphRAG v1 perform no embedding, so the setting does not apply to them.

\paragraph{Comparability.} The prefixed 100-question phone runs contain all 20 questions of the
primary phone comparison. Apart from the prefixes, they use the same reader, retrieval depth,
history budget, prompt, decoding and judge. On those 20 questions, the prefixes changed Gemma~4 E2B
from 10 to 8 correct answers and E4B from 8 to 10. Each reader had two discordant questions, in
opposite directions (exact McNemar $p=0.5$ for each). At this sample size, the phone shows no
detectable effect of the prefixes on accuracy. The check covers accuracy only: the 100-question
sittings spread their memory build over five times as many questions, so their energy per answer is
not comparable. Workstation prefix gains are not applied to phone scores, and every phone
efficiency result uses answers and energy from the same measured workload.

\section{Answer scoring and judge characterization}
\label{sec:judge}

Answer correctness is decided by an automatic judge: Llama~3.1 70B Instruct at \texttt{Q4\_K\_M}
precision~\citep{llama31}, run on the workstation. The judge is prompted with the official
LongMemEval correctness rubric, which supplies the question, the reference answer, and the model's
prediction, and asks whether the response contains or is equivalent to the reference answer. This
correctness prompt is distinct from the answer prompt given to the reader. LoCoMo answers use its
standard form; LongMemEval answers use the benchmark's own variants for temporal, knowledge-update,
preference and abstention questions (Appendix~\ref{app:repro_prompts}). Answers produced on the phone are
transferred to the workstation and scored by the same judge under the same rubric. Every accuracy
figure in the thesis, on either device, comes from this judge and is stamped with the judge's
identity so that its provenance is traceable.

\paragraph{Judge characterization and reliability.} The judge is treated as an instrument with
measurable noise, not an infallible oracle. Its consistency was checked by re-judging identical
prediction sets at $T=0$. Three back-to-back re-judgings of 500 predictions in one server session
changed no verdict. When 100 predictions were judged again in a later session, one verdict changed;
GPU-level non-determinism is a plausible cause, and a single comparison of this size does not give a
tight bound. Lexical measures (exact match, token overlap, substring containment) are computed
as sanity checks, to confirm that the relative ranking of the strategies is not an artifact of judge
leniency.

\paragraph{Human validation.} Because an automatic judge cannot validate itself, a blind human audit
checked its verdicts on 100 judged answers. The answers were drawn from eight runs covering both
benchmarks, all three memory strategies and both Gemma readers, about half of them judged correct and
half judged incorrect; some come from runs made before the embedding-prefix correction, which does not
affect what the judge is asked to decide. Every condition identifier was removed and the order was
shuffled. Because the runs answer the same questions, the 100 answers cover 86 distinct questions.
The author labelled each answer as correct, incorrect or ambiguous against the reference answer, using
the judge's rubric, without seeing the judge's verdicts.

On the 83 answers given a decisive label, the judge agreed with the human label on 82 (98.8\%, Cohen's
$\kappa=0.976$); keeping one answer per question gives 73 of 74. The single disagreement was a
false positive: the reference was ``fitness tracker'' and the judge accepted the hedged answer
``fitness tracker or smartwatch''. The judge marked no correct answer as wrong. The other 17 answers
were labelled ambiguous. Thirteen were LoCoMo answers that matched the reference only partly or
loosely, such as a specific date where the reference names a week; four were LongMemEval answers,
three of them to abstention questions. Ambiguous answers were not assigned a label and are left out
of the agreement denominator, which is therefore 83. The judge had accepted 14 of the 17. Counting
every ambiguous answer as a disagreement instead would lower agreement to 82 of 100. The audit has one annotator,
so it does not measure agreement between humans, and no second judge model was compared. Because
the sample is balanced between judged-correct and judged-incorrect answers, it measures how often
the judge's verdicts are right, not how far any particular accuracy figure would move.

\paragraph{Missing data and error handling.} Model generation failures and judge infrastructure
errors are kept distinct. Predictions that are genuinely empty, contain only unclosed formatting
tags, or exhaust the generation budget before producing an answer are kept in the denominator and
scored as incorrect. Transient connection failures or timeouts during judging are retried with
exponential backoff; verdicts that still cannot be parsed are isolated and logged rather than counted
as wrong.

\section{On-device resource measurement}
\label{sec:energy}

Physical resource measurements are made on the Samsung Galaxy S23 Ultra.

\paragraph{Measurement environment.} Runs are performed with the device unplugged, held under a
continuous wake-lock so the processor does not sleep, and with the display off.\footnote{The 22
sittings behind the main phone tables started at 86 to 100\% charge and a battery temperature of
31.7 to $36.8^\circ\text{C}$, and each reached 38.6 to $43.3^\circ\text{C}$ within ten minutes.
Eight of them, and three further runs, started at the highest possible charge (100\%), where the
phone's charge counter can pause for the first few minutes, so those runs under-count their energy
(Appendix~\ref{app:ch5_energy_conventions}).}

\paragraph{Energy method of record.} Battery-side energy is estimated from the phone's charge counter
and measured voltage. The counter's charge decrease agrees with the reported battery-level decrease
to within about one percentage point. A separate cross-check using integrated current and voltage
agrees with counter-based energy to within 10\% on sittings longer than 50 minutes
(Appendix~\ref{app:ch5_energy_conventions}). Resource telemetry is sampled every 0.5~seconds by a background
profiler logging directly to file. The physical quantity directly measured on the handset is the battery
charge consumed, $\Delta Q$ (in $\text{mAh}$), recorded by the battery's hardware Coulomb counter:
\begin{equation}
  \label{eq:charge_formula}
  \Delta Q = \frac{Q(t_0)-Q(t_N)}{1000}.
\end{equation}
Here, $Q(t)$ is the charge counter in $\mu\text{Ah}$ and $t_0,\ldots,t_N$ are the usable unplugged
samples inside the measured run span. Whole-device battery energy is then calculated by pairing each
charge interval with the mean battery voltage measured over that interval:
\begin{equation}
  \label{eq:energy_formula}
  E = 3.6 \times 10^{-3}
      \sum_{i=0}^{N-1}
      \left[Q(t_i)-Q(t_{i+1})\right]
      \frac{V(t_i)+V(t_{i+1})}{2},
\end{equation}
where $V(t)$ is measured in volts and $E$ is obtained in joules. The factor
$3.6\times10^{-3}$ converts $\mu\text{Ah}$ to coulombs. This voltage-weighted Coulomb-counting
calculation uses the phone's recorded voltage, not a fixed nominal-voltage assumption. The
study reports both $\text{mAh}$ and $\text{J}$: the former is the direct charge measurement, while
the latter is the corresponding battery-side energy estimate.

\paragraph{On-device efficiency metric.} Accuracy $a$ measures the fraction of answers judged correct,
while energy per query, $E/N$, measures the cost of producing an answer, whether correct or wrong.
Both are reported separately. Their combined measure asks how many correct answers the workload
produces per unit of energy. For a run containing $N$ scored answers with phone-measured accuracy
$a$, correct answers per unit of energy are
\begin{equation}
  \label{eq:golden_metric}
  G_E = \frac{aN}{E/1000},
\end{equation}
reported per kilojoule. The denominator includes the energy consumed over the recorded workload,
including queries answered incorrectly. Answer correctness does not validate the energy instrument;
the telemetry cross-checks and instrument limitations are assessed separately.
When a full-charge equivalent is reported, it is calculated
directly from charge use as $G_Q = aN \times 5000 / \Delta Q$, where $5000\,\text{mAh}$ is the rated
battery capacity. Thus the energy-normalized metric uses the measured voltage trace, whereas the
full-charge metric does not require an assumed battery voltage.

\paragraph{Instrument limits.} Energy uses the hardware Coulomb counter, which integrates charge
continuously at the battery controller, rather than polled instantaneous current, which drifts under
rapid load changes; the current reading serves only as a cross-check. The counter moves in 0.1\% steps
and holds its value at full charge, which affected nine of the eleven sittings that started at 100\%
(Appendix~\ref{app:ch5_energy_conventions}).

\paragraph{Physical telemetry and resource metrics.} Beyond battery charge, the background profiler
requests a sample every 0.5~seconds. Each battery query through Termux:API takes a second or more,
so the realized interval was 1.2 to 3.2~seconds (median per run). Each sample records:
\begin{itemize}
  \item \textbf{Battery charge and energy.} Directly measured charge drop ($\Delta Q$ in $\text{mAh}$)
        and derived whole-device energy estimate ($E$ in Joules) consumed across the entire smartphone
        (SoC, memory, display-off circuitry, background services) from initial ingestion to the final
        emitted token.
  \item \textbf{CPU clock frequencies.} Active clock frequencies (in MHz) sampled across
        all eight CPU cores (the Cortex-X3 prime core and the Cortex-A7x and A510 clusters) via
        \texttt{/sys/devices/system/cpu/cpu*/cpufreq/} to monitor frequency scaling. A frequency
        trace shows when cores slowed, but not whether the thermal governor caused it.
  \item \textbf{Temperature.} The hottest reading across all thermal zones that can be read without
        root (per-core CPU, GPU and other sensors), in degrees Celsius ($^\circ\text{C}$). It is a
        device-wide maximum, not one SoC sensor, and it is not the battery temperature, which
        is logged separately. Sampling starts with the harness, after the model servers have loaded,
        so the first sample is not an idle baseline.
  \item \textbf{Memory footprint (RAM).} System used and available physical memory (in MB) sampled via
        \texttt{/proc/meminfo} to evaluate memory pressure and KV-cache overhead during long-context
        processing.
  \item \textbf{Latency.} End-to-end wall-clock time per question, with the prompt and answer token
        counts reported by the server. The phone records only the total generation time, so prefill
        and decode time are not separated.
\end{itemize}

\paragraph{Replicates and reporting conventions.} Vector retrieval on all five readers and full
context on Gemma E2B were each run in three independent sittings, to measure how
much the energy varies between runs. Full context on Gemma E4B and on the three Qwen readers used
one sitting each. Because a fixed-question run spreads its one-time memory-construction cost over
only twenty questions, per-answer energy includes building the memory, and the ratios between
strategies are specific to this workload. Graph-construction energy on the phone comes from one
directly measured 419-turn build ($21.9\,\text{kJ}$, or $52.3\,\text{J/turn}$); any value at other
conversation lengths is a projection from a workstation-measured slope and is labelled as such.

\section{Controls, validity, and statistical plan}
\label{sec:validity}

The statistical analysis has three specific purposes. It compares memory strategies on the same
questions, checks whether the grouping of LoCoMo questions by conversation affects those comparisons,
and tests whether workstation accuracy can be used in place of phone accuracy. The three procedures
below are kept separate because they answer different questions.

\paragraph{Paired strategy comparisons.} For a fixed benchmark and reader, every memory strategy
answers the same questions. Each answer is reduced to a binary outcome (correct or incorrect), and
the strategies are compared with McNemar's exact test~\citep{mcnemar}. The test uses only the
discordant pairs: questions answered correctly by one strategy and incorrectly by the other. Results
report the accuracy of each strategy, the difference in percentage points, the two discordant counts,
and the exact $p$-value. Statistical claims are limited to the main comparisons between full context,
vector retrieval, and the graph versions. Retrieval-depth sweeps, prompt experiments, and other
ablations are treated as descriptive analyses rather than separate significance claims.

\paragraph{Conversation-level sensitivity.} The 200 LoCoMo questions come from ten multi-session
conversations, so questions from the same conversation may be correlated. A conversation-level block
bootstrap is used as a sensitivity analysis: complete conversations are resampled with replacement,
and the paired accuracy difference is recalculated for each bootstrap sample. With only ten
conversation clusters, this analysis checks the stability of the result rather than replacing the
question-level McNemar test. LongMemEval-S contains 60 questions drawn from 60 separate sessions and
therefore does not use this conversation-level check.

\paragraph{Cross-environment agreement check.} The workstation and phone use different runtimes and
weight packages, so workstation accuracy is not assumed to transfer to the phone. The same 100 E2B
vector-retrieval questions were evaluated in both environments. The comparison reports the percentage
of matching correctness labels, Cohen's $\kappa$, and an exact McNemar test for directional bias. The
decision rule was fixed before the comparison: workstation accuracy could be substituted only if
agreement reached at least 90\% and McNemar's test gave $p>0.05$. The outcome is reported in
Section~\ref{sec:ch5_parity}.

\paragraph{Answer-prompt robustness.} The reader is given one curated short-answer prompt across all
conditions and both benchmarks, so that strategy and benchmark comparisons are not confounded with
prompt wording. Because a fixed prompt could still interact with the context representation, the
LoCoMo grid was re-run with LoCoMo's own official question prompt, an externally defined standard, in
place of the curated one, with everything else held fixed (Section~\ref{sec:ch4_robustness}). The
substitution was not evaluated on LongMemEval-S, which limits claims about absolute accuracy on that
benchmark.

\paragraph{Reporting conventions.} Every result in Chapters~\ref{ch:memory_evaluation}
to~\ref{ch:spe} carries one of four labels. Tables~\ref{tab:ch9_claims} and~\ref{tab:ch9_claims_b} list the main claims with
their label, sample size and test.
\begin{itemize}
  \item \emph{Pre-registered}: a criterion fixed before the run that tests it. Only the parity
        threshold above and the two SPE criteria of Chapter~\ref{ch:spe} have this status, and
        only these are confirmatory tests.
  \item \emph{Primary}: nine paired comparisons, eight on Gemma and one on Qwen3.5 4B, reported with
        their exact McNemar $p$-value and a Holm correction across the nine. This family was
        designated during the analysis, after the results were available, so it is not
        pre-registered.
  \item \emph{Exploratory}: every other paired test, reported with its unadjusted exact $p$-value.
  \item \emph{Descriptive}: measurements reported without a test, including all phone energy
        and efficiency ratios.
\end{itemize}
Four readings apply throughout and are not repeated in each chapter. A difference that does not
reach significance is never read as evidence that two conditions are equivalent. Workstation
results carry no phone-energy measurement, and workstation accuracy never stands in for phone
accuracy. Graph conclusions concern graphs extracted by the tested small models and retrieved by
turn-anchored triple selection, not knowledge graphs in general. Finally, restoring source text
adds both content and prompt length, so its effect does not isolate the information lost when
dialogue is reduced to triples.

\section{Software and reproducibility}
\label{sec:repro}

The pipeline is implemented in Python and coordinates a llama.cpp inference server on the phone
(under Android Termux) and Ollama on the Linux workstation. The graph is implemented with
NetworkX~\citep{networkx}, and embeddings use Nomic Embed Text v1.5~\citep{nomic_embed}.

Every result, table, and statistical test reported in the thesis is produced by scripts directly from
the stored per-run logs rather than copied by hand: one script aggregates the master accuracy table,
one derives the on-device energy metric and is the sole method of record for that number, and one
produces the statistical report. Hand-computed intermediate figures from working sessions are not
used in the text.

To keep the exact artifacts identifiable, an artifact manifest collects the runtime build versions,
model repository revisions and file hashes, and the exact prompt and template texts used in the
measured runs. The model revisions and hashes were reconstructed afterwards from the surviving model
cache and server logs, and the build commit of the phone runtime could not be recovered
(Section~\ref{sec:ch5_boundary}). The exact identifiers, command flags, and prompt texts are listed
in the reproducibility appendix.
